\documentclass{article}
\pagestyle{empty}

%% >>> preamble
\usepackage{amsmath}
\usepackage{amsthm}
\theoremstyle{plain}
\newtheorem{theorem}{Theorem}[section]
\newtheorem{lemma}[theorem]{Lemma}
\newtheorem{proposition}[theorem]{Proposition}
\theoremstyle{definition}
\newtheorem{definition}[theorem]{Definition}
\theoremstyle{remark}
\newtheorem*{remark}{Remark}
%% <<< preamble

\begin{document}
%% >>> snippet
\section{Convexity}
\begin{definition}[Convex function]
  A function $f$ is \emph{convex} if
  $f(\lambda x + (1-\lambda) y) \le \lambda f(x) + (1-\lambda) f(y)$
  for all $x, y$ and all $\lambda \in [0, 1]$.
\end{definition}

\begin{theorem}\label{thm:global}
  Every local minimum of a convex function is a global minimum.
\end{theorem}

\begin{proof}
  Suppose $x^\star$ is a local but not a global minimum, so some $y$ has
  $f(y) < f(x^\star)$. Convexity gives points arbitrarily close to $x^\star$
  with smaller value, a contradiction.
\end{proof}

\begin{remark}
  Strict convexity also makes the minimiser unique.
\end{remark}
%% <<< snippet
\end{document}
